\Article{34}
It would be foreign to our present purpose to pass from the general principles established here to particular phenomena, especially since the substance of those principles agrees with the theory by which the illustrious \textsc{Laplace} has already explained a great many phenomena observed in the equilibrium of fluids with equal skill and success. A wide field certainly remains, promising a rich new harvest, but let us leave it for future work. It will be appropriate, however, to add a few remarks that may either shed new light on the subject or guard against a mistaken interpretation.

I. Our theory makes no claim to determine the equilibrium shape with mathematical exactness. It is content to determine a shape from which the true equilibrium shape cannot differ by a perceptible amount. You would be mistaken to attribute this to any imperfection of the theory. It has accomplished in full what can be accomplished as long as the law of molecular attraction is unknown. In the state of equilibrium the function $\Omega$ must be exactly at a maximum, and hence the function
\[
\frac{2\pi cs\psi0}g-\frac\Omega{gc}
\]
at a minimum. Owing to the nature of molecular attraction, this function is only imperceptibly different from the function $W$, though it is not exactly equal to it. Thus the shape for which $W$ is a minimum is not the exact equilibrium shape, but the difference must be imperceptible insofar as any perceptible change of that shape would produce a perceptibly greater value of the function $W$. This clearly does not rule out a perceptible difference in the curvature of the surface, provided it is confined to an imperceptible part of the surface. Accordingly, in the exact equilibrium shape the constant angle denoted above by $A$ can be regarded as the inclination of the surface of the fluid to the vessel wall only at an immeasurable distance from the vessel and no longer at the point of contact itself. Or, as the illustrious \textsc{Laplace} rightly noted, the inclination at the boundary of the sphere of perceptible attraction of the vessel agrees perceptibly with the value of $A$.

\SourceLeaf{79}II. We must carefully distinguish the equilibrium shape from the shape at rest. Whenever a fluid is in equilibrium, it must certainly remain in that state. Yet whenever the shape of the fluid differs somewhat from the equilibrium shape, the fluid may nevertheless remain at rest or, if it is moving, may already lose its motion before it reaches equilibrium. In the same way, for example, a cube simply placed on a horizontal plane is in equilibrium, but it can also rest on an inclined plane when friction prevents motion. Thus a fluid in a state for which $W$ has a minimum value will certainly be at rest. But whenever the fluid is in a different state, for instance when $W$ can decrease, it will pass from that state to equilibrium only insofar as friction does not prevent it. In this respect the two conditions of equilibrium differ essentially. The first fundamental equation (Article 28) is independent of the ability of the boundary $P$ to move, that is, it is necessary for the minimum condition even when the boundary is assumed fixed. Thus insofar as the fluid has perfect fluidity, so that one part can slide freely over another when even the smallest force calls for motion, the fluid will necessarily conform to that condition. The second principle (Article 29) is quite different, for it depends essentially on the complete freedom of the boundary $P$ to move on the vessel’s surface. The condition for the value of $W$ to be a minimum does require the equation $i=A$. If, however, after the fluid surface has conformed to the first principle, the angle $i$ has yet to reach its normal value, and thus $W$ has yet to attain its absolute minimum, the transition to perfect equilibrium requires a shift of the boundary $P$, or equivalently motion of the fluid in contact with the vessel, and friction can indeed oppose such motion. This makes it clear why we encounter such great differences in the value of the angle $i$ in experiments on the same bodies. Likewise, when $\beta\beta>\alpha\alpha$, a fluid in a vessel whose walls are already wetted will indeed settle according to the law of equilibrium, under which the substantial part of the fluid must have $i=180^0$. But in a vessel whose walls are still dry beyond the fluid, a fluid starting from a state out of equilibrium and spreading over the dry vessel walls can come to rest before the angle $i$ has reached the value $180^0$. This also explains why the capillary phenomena of fluids that do not resist wetting show such irregularities in dry tubes, with the rise very often much smaller than in tubes already wetted, where we always find wonderfully close agreement with the theory.

\expandafter\def\csname HistoricalNoteText7\endcsname{It is worth noting that the quantity denoted by the illustrious \textsc{Laplace} by $H$ agrees with our $\pi c\theta0$, so that his $a$ denotes what in our notation is $\frac{g}{\pi c\theta0}$ or $\frac1{2\alpha\alpha}$.}%
\SourceLeaf{80}III. The ratio of the constants $\alpha,\beta$ cannot be derived from the phenomena whenever $\beta$ is greater than $\alpha$. In this case the shape of the same fluid in vessels of equal shape but different materials differs only in respect of the immeasurable film wetting the vessel. Whenever $\beta$ is smaller than $\alpha$, the ratio of these constants can indeed be determined by means of the angle $i$, but for the reasons just given the determination can hardly be very precise. For mercury in glass vessels the illustrious \textsc{Laplace} set the angle $i=43^0\,12'$.

The constant $\alpha$ can be determined with much greater precision, especially if we can use vessels that can be wetted. For water at a temperature of $8{,}5$ degrees on the centigrade thermometer we must set, according to the experiments cited by the illustrious \textsc{Laplace}\HistoricalNote{7},
\[
\begin{aligned}
\alpha\alpha&=7{,}5675\ \text{square millimetres, or}\\
\alpha&=2{,}7509\ \text{millimetres}
\end{aligned}
\]
For spirit of wine, whose specific gravity is $=0{,}81961$, at the same temperature,
\[
\begin{aligned}
\alpha\alpha&=3{,}0441\ \text{square millimetres, or}\\
\alpha&=1{,}7447\ \text{millimetres}
\end{aligned}
\]
For oil of turpentine at a temperature of $8$ degrees,
\[
\begin{aligned}
\alpha\alpha&=3{,}305\ \text{square millimetres, or}\\
\alpha&=1{,}818\ \text{millimetres}
\end{aligned}
\]
For mercury at a temperature of $10$ degrees, until new experiments afford greater precision, we may set
\[
\begin{aligned}
\alpha\alpha&=3{,}25\ \text{square millimetres, or}\\
\alpha&=1{,}803\ \text{millimetres}
\end{aligned}
\]
It is likely, moreover, that temperature affects the value of the constant $\alpha\alpha$ only insofar as density depends on it. On this hypothesis the value of $\alpha\alpha$ will therefore be proportional to the density.

\SourceLeaf{81}These values were inferred from the rise or depression of fluids in capillary tubes. It is very difficult, however, to measure the diameters of those tubes exactly, and more difficult still to establish that their cross sections are circular. Much greater precision is promised by experiments on the diameters and volumes of large drops of mercury resting on a horizontal plate or on one of very slight known curvature, such as those already conducted by the physicists \textsc{Segner} and \textsc{Gay-Lussac}. So too, for liquids that wet glass vessels, are experiments on the dimensions of large air bubbles in vessels closed at the top with a wetted cover that is flat and horizontal or slightly curved to a known radius. We invite physicists to undertake these experiments.

IV. To avoid exceeding the limits set for this paper, we have had to restrict the application of our general principles here to the simplest case, in which we consider a single liquid in a rigid vessel. Nothing prevents us, however, from giving the theory its fullest generality, so that it also embraces the problem of several liquids in the same vessel, as well as the case in which rigid bodies, either wholly or partly free, are also immersed in the fluid. But we must reserve a fuller treatment of these questions for another occasion.
